\subsection*{数据搬运与分块选择}

下面汇总数据搬运和分块选择的常用写法。

\begin{table}[htbp]\centering\small
\begin{tabularx}{\linewidth}{@{}lX@{}}
\toprule
语法 & 含义 \\
\midrule
\texttt{dma.copy.async ...} & 启动非阻塞 DMA 传输。 \\
\texttt{wait f} & 等待 \texttt{f} 对应的传输完成。 \\
\texttt{future.data} & 访问已搬运分块。 \\
\texttt{future.span} & 访问已搬运分块的形状。 \\
\texttt{dma.transp\textless{}...\textgreater{}} & 搬运过程中置换维度。 \\
\texttt{dma.pad\textless{}...\textgreater{}} & 搬运过程中填充，并显式给出填充值。 \\
\texttt{.zfill} & 把不完整分块中的越界位置补零。 \\
\texttt{.swiz\textless{}N\textgreater{}} & 选择共享内存重排布局。 \\
\texttt{chunkat(...)} & 从规则分块中选出一个等分块。 \\
\texttt{subspan(...).at(...)} & 按显式分块范围与分块坐标选取分块。 \\
\texttt{view(...).from(...)} & 按显式范围与显式元素偏移选取分块。 \\
\texttt{tile \# i} & 把分块坐标与块内坐标合成为全局坐标。 \\
\texttt{\#tile} & 求 \texttt{tile} 轴的范围。 \\
\bottomrule
\end{tabularx}
\end{table}

阅读搬运语句时，先确定所选区域和目标位置，再检查目标何时就绪。流水线中仍需回答这两个问题；额外缓冲区只是让不同分块能够同时处在这个过程的不同阶段。
